% Modelo da UFMG -
% Este modelo foi baseado em: modelo-ufpr.tex,v 1.1 2003/06/30 15:05:18 gweber Exp $
% $Id: modelo-ufpr.tex,v 1.1 2003/06/30 15:05:18 gweber Exp $
%   Licence: LPPL (LaTeX Project Public License)
%     You can change this file in the terms of LPPL
%     (http://www.latex-project.org/lppl.html)
% copyright Rogério C. <rogerioc@cesec.ufpr.br>
%
% ****** DEFINIÇÕES INICIAIS ******	
\documentclass[a4paper,11pt, normaltoc, pnumromarab, pagestart=fistchapter, tocpage=plain]{abnt}
% Utilize a opção normalfigtabnum para numerar as figuras e tabelas por capítulo
%\usepackage[alf]{abntcite} chamada de referencia alfabetica

\usepackage[brazil]{babel}
\usepackage[utf8]{inputenc}
\usepackage[T1]{fontenc}
\usepackage{indentfirst}
\usepackage{graphicx}														%Package para figuras
\usepackage{float}
\usepackage{geometry}
\usepackage{fancyhdr}
\pagestyle{fancy}
\usepackage{url}

%\geometry{a4paper,left=3cm,right=3cm,top=3cm,bottom=3cm}

%%
%%	Ainda em teste
%%
%\usepackage[bookmarks=false]{hyperref}					%Package para hyper-referências
%\hypersetup{colorlinks,
%							citecolor = black,
%							filecolor = black,
%							linkcolor = black,
%							urlcolor  = blue,
%						pdfnewwindow}
%
% O problema ocorre quando há referências do tipo \cite{} e \citeonline{}
% Há ainda outros problemas -> o figura, antes do número, não altera de cor na lista de figuras.
% O mesmo ocorre na lista de tabelas.
% O sumário aponta para a capa, não para o resumo, lista, apêndices ou anexos correspondentes.
% -> Funciona para capítulos.
%

\makeatletter	%Para que ele entenda o @

% Altera o tamanho das fontes dos capítulos e dos apêndices
\renewcommand{\ABNTchapterfont}{\bfseries}
\renewcommand{\ABNTchaptersize}{\Large}
\renewcommand{\ABNTanapsize}{\Large}

%Altera o espaçamento entre dots
%\renewcommand\@dotsep{2}




%Altera forma de montagem do table of contents
\renewcommand\l@chapter[2]{
  \ifnum \c@tocdepth >\m@ne
    \addpenalty{-\@highpenalty}%
    \vskip 1.0em \@plus\p@
    \setlength\@tempdima{1.5em}%
    \begingroup
      \ifthenelse{\boolean{ABNTpagenumstyle}}
        {\renewcommand{\@pnumwidth}{3.5em}}
        {}
      \parindent \z@ \rightskip \@pnumwidth
      \parfillskip -\@pnumwidth
      \leavevmode \normalsize\ABNTtocchapterfont
      \advance\leftskip\@tempdima
      \hskip -\leftskip
      #1\nobreak\dotfill \nobreak%
      \ifthenelse{\boolean{ABNTpagenumstyle}}
         {%
          \hb@xt@\@pnumwidth{\hss
            \ifthenelse{\not\equal{#2}{}}{{\normalfont p.\thinspace#2}}{}}\par
         }
         {%
          \hb@xt@\@pnumwidth{\hss #2}\par
         }
      \penalty\@highpenalty
    \endgroup
  \fi}

\renewcommand*\l@section{\@dottedtocline{1}{0em}{2.3em}}
\renewcommand*\l@subsection{\@dottedtocline{2}{0em}{3.2em}}
\renewcommand*\l@subsubsection{\@dottedtocline{3}{0em}{4.1em}}

% Cria um comando auxiliar para montagem da lista de figuras
\newcommand{\figfillnum}[1]{%
  {\hspace{1em}\normalfont\dotfill}\nobreak
  \hb@xt@\@pnumwidth{\hfil\normalfont #1}{}\par}

% Cria um comando auxiliar para montagem da lista de tabelas
\newcommand{\tabfillnum}[1]{%
	{\hspace{1em}\normalfont\dotfill}\nobreak
	\hb@xt@\@pnumwidth{\hfil\normalfont #1}{}\par}

% Altera a forma de montagem da lista de figuras
\renewcommand*{\l@figure}[2]{
	\leftskip 3.1em
	\rightskip 1.6em
	\parfillskip -\rightskip
	\parindent 0em
	\@tempdima 2.0em
	\advance\leftskip \@tempdima \null\nobreak\hskip -\leftskip
	{Figura \normalfont #1}\nobreak \figfillnum{#2}}

% Altera a forma de montagem de lista de tabelas
\renewcommand*{\l@table}[2]{
	\leftskip 3.4em
	\rightskip 1.6em
	\parfillskip -\rightskip
	\parindent 0em
	\@tempdima 2.0em
	\advance\leftskip \@tempdima \null\nobreak\hskip -\leftskip
	{Tabela \normalfont #1}\nobreak \tabfillnum{#2}}

% Define os comandos que montam a lista de símbolos
\newcommand{\listadesimbolos}{\pretextualchapter{Lista de Símbolos}\@starttoc{lsb}}
\newcommand{\simbolo}[2]{{\addcontentsline{lsb}{simbolo}{\numberline{#1}{#2}}}#1}
\newcommand{\l@simbolo}[2]{
	\vspace{-0.75cm}
	\leftskip 0em
	\parindent 0em
	\@tempdima 5em
	\advance\leftskip \@tempdima \null\nobreak\hskip -\leftskip
	{\normalfont #1}\hfil\nobreak\par}

% Define o comando que monta a lista de siglas
\newcommand{\listadesiglas}{\pretextualchapter{Lista de Siglas}\@starttoc{lsg}}
\newcommand{\sigla}[2]{{\addcontentsline{lsg}{sigla}{\numberline{#1}{#2}}}#1}
\newcommand{\l@sigla}[2]{
	\vspace{-0.75cm}
	\leftskip 0em
	\parindent 0em
	\@tempdima 5em
	\advance\leftskip \@tempdima \null\nobreak\hskip -\leftskip
	{\normalfont #1}\hfil\nobreak\par}

% Define o tipo de numeração das páginas
\renewcommand{\chaptertitlepagestyle}{plain}

% Altera a posição da numeração de páginas dos elementos pré-textuais
\renewcommand\pretextualchapter{
	\if@openright\cleardoublepage\else\clearpage\fi
	\pagestyle{\chaptertitlepagestyle}
	\global\@topnum\z@
	\@afterindentfalse
	\@schapter}

% Altera a posição da numeração de páginas dos elementos textuais
\renewcommand{\ABNTchaptermark}[1]{
	\ifthenelse{\boolean{ABNTNextOutOfTOC}}
		{\markboth{\ABNTnextmark}{\ABNTnextmark}}
		{\chaptermark{#1}
}}


% Redefine o tipo de numeração das páginas
\renewcommand{\ABNTBeginOfTextualPart}{
	\renewcommand{\chaptertitlepagestyle}{plain}
	\renewcommand{\thepage}{\arabic{page}}
%	\setcounter{page}{1}
}

% set headings
\lhead{}

\rhead
      {\fancyplain{}{\footnotesize\rightmark}}

\cfoot{}

\makeatother

%Altera o tamanho do parágrafo
\setlength{\parindent}{1.5cm}

% ********************************
% ***** Início do Documento ******
% ********************************
\begin{document}


% set headings


%    Modelo de Projeto de TCC
%
% Centro Federal de Educação Tecnológica de Minas Gerais - CEFET-MG
% Autor: Mariane Raquel
%
% Parte: Capa e Folha de rosto projeto

\titulo{Comparação de algoritmos paralelos de ordenação em MapReduce}

\instituicao{Centro Federal de Educação Tecnológica de Minas Gerais \par 					Departamento de Computação \par
			Curso de Engenharia de Computação}

\autor{Mariane Raquel Silva Gonçalves}

\orientador{Profª. Drª. Cristina Duarte Murta}

\comentario{Projeto apresentado como requisito parcial à aprovação em Trabalho de Conclusão de Curso I do curso de Engenharia de Computação do Centro Federal de Educação Tecnológica de Minas Gerais.}

\local{Belo Horizonte}

\data{2012}

\logo{figuras/logo.pdf}

\capa

\folhaderosto


%% -------------------------------------------------------
%% --------------------- CAPA ----------------------------
%% -------------------------------------------------------
%
%\begin{titlepage}
%
%% Logo e Universidade
%\begin{flushleft} 
%\begin{minipage}{0.25\textwidth}
%\includegraphics[scale=0.7]{figuras/logo.png} 
%\end{minipage} \hfill
%\begin{minipage}{0.7\textwidth}
%\begin{flushleft} \begin{center} 
%Centro Federal de Educação Tecnológica de Minas Gerais \\ Departamento de Computação \\ Curso de Engenharia de Computação 
%\end{center} \end{flushleft} 
%\end{minipage} 
%\end{flushleft}
%\begin{center}
%% Aluno 
%\vspace{1.5cm} Mariane Raquel Silva Gonçalves
%% Título
%\vspace{5cm} 
%\begin{espacoduplo} 
%\begin{Large} \textbf{ \textsc{Comparação de algoritmos paralelos de ordenação em MapReduce}}  \end{Large} \\
%\end{espacoduplo}
%% Orientador
%\vfill Orientadora: Profª. Drª. Cristina Duarte Murta \\[1cm]
%% Data
%Belo Horizonte \\ 2012 \\
%\end{center}
%\end{titlepage}
%
%
%% -------------------------------------------------------
%% ----------------- FOLHA DE ROSTO ----------------------
%% -------------------------------------------------------
%\instituicao{Centro Federal de Educação Tecnológica de Minas Gerais \\ Departamento de Computação \\ Curso de Engenharia de Computação}
%
%\autor{ Mariane Raquel Silva Gonçalves} 
%
%\titulo{Comparação de algoritmos paralelos de ordenação em MapReduce}
%
%%\tituloformat = sc, bf, Large
%
%\orientador{Profª. Drª. Cristina Duarte Murta} 
%
%%LOGO \includegraphics[scale=0.7]{figuras/logo.png}
%
%\comentario{Projeto apresentado como requisito parcial à aprovação em Trabalho de Conclusão de Curso I do curso de Engenharia de Computação do CEFET-MG. }
%
%\folhaderosto
%%\begin{titlepage}
%%
%%% Logo e Universidade
%%\begin{flushleft} 
%%\begin{minipage}{0.25\textwidth}
%%\includegraphics[scale=0.7]{figuras/logo.png} 
%%\end{minipage}
%%\hfill
%%\begin{minipage}{0.7\textwidth}
%%\begin{flushleft} 
%%\begin{center} Centro Federal de Educação Tecnológica de Minas Gerais \\ Departamento de Computação \\ Engenharia de Computação \end{center}
%%\end{flushleft}
%%\end{minipage}
%%\end{flushleft}
%%
%%\begin{center}
%%% Aluno 
%%\vspace{1.5cm} Mariane Raquel Silva Gonçalves
%%% Título
%%\vspace{5cm} 
%%\begin{espacoduplo}  \begin{Large} \textbf{ \textsc{Comparação de algoritmos paralelos de ordenação em MapReduce}}  \end{Large} \\ \end{espacoduplo}
%%\end{center}
%%
%%% Texto
%%\vfill
%%\begin{flushright}
%%\begin{minipage}{0.5\textwidth}
%%Projeto apresentado como requisito parcial à aprovação em Trabalho de Conclusão de Curso I do curso de Engenharia de Computação do CEFET-MG. 
%%\end{minipage}
%%\end{flushright}
%%\vfill
%%
%%\begin{center}
%%% Orientador
%%Orientadora: Profª. Drª. Cristina Duarte Murta \\[1cm]
%%% Data
%%Belo Horizonte \\ 2012 \\ 
%%\end{center}
%%\end{titlepage}
%



\sumario

% 1 - Lista de Figuras
%\listadefiguras

% 2 - Lista de Tabelas
%\listadetabelas

% 3 - Lista de Siglas
% Forma de uso: \sigla{sigla}{Descrição}
%\listadesiglas


\chapter{Proposta de Projeto}

Este texto apresenta uma proposta de projeto para a disciplina Trabalho de Conclusão de Curso. 

\section{Tema do projeto}

[Alguma citação que diga algo assim: com o surgimento das arquiteturas multicore, a computação paralela é hoje o caminho para desenvolver sistemas. Possivelmente Tally 2007, Sutter e Larus 2005 => sugiro os artigos [Asanovic et al. 2009]ASANOVIC, K. et al. A view of the parallel computing landscape.
Commun. ACM   e [Kale e Solomonik 2010] e mais alguns, veja em anexo] .

Uso crescente de computação paralela para sistemas computacionais gera a necessidade de algoritmos de ordenação inovadores, desenvolvidos para dar suporte a essas aplicações. A ordenação paralela é o processo de reorganizar uma sequência de entrada e produzir uma saída ordenada de acordo com um atributo através de múltiplas unidades de processamento, que ordenam em conjunto a sequência de entrada. 

Na criação de algoritmos de ordenação paralela, é ponto fundamental ordenar coletivamente os dados de cada processo individual, de forma a utilizar todas as unidades de processamento e minimizar os custos de redistribuição de chaves entre os processadores \cite{Kale:2010}. 

\section{Relevância}

%1. porque computação paralela

O limite físico de aumento na frequência de operação dos processadores levou a indústria de hardware a substituir o processador de núcleo único por vários processadores eficientes em um mesmo \textit{chip}, os processadores \textit{multicore}. 
É preciso, então, criar aplicações que utilizem efetivamente o crescente número de núcleos dos processadores  \cite{Asanovic:2009}.

%Com a recente mudança de foco no desenvolvimento de algoritmos sequenciais para paralelos, conhecer os modelos de programação paralela se tornou uma grande necessidade na computação. 
%. No entanto,  Os algoritmos paralelos ainda são um ramo pouco explorado, devido a  maior complexidade no desenvolvimento e recentes aplicações em sistemas \textit{multicore}. 

%2. porque ordenação paralela

Um grande número de aplicações paralelas possui uma fase de computação intensa, na qual um lista de elementos deve ser ordenada com base em algum de seus atributos. Um exemplo é o algoritmo de Page Rank \cite{PageRank:1999} da Google: as páginas de resultado de uma consulta são rankeadas de acordo com sua relevância, e então precisam ser ordenadas de maneira eficiente \cite{Kale:2010}.

Fatores como movimentação de dados, balanço de carga, latência de comunicação e distribuição inicial das chaves são considerados ingredientes chave para o bom desempenho da ordenação paralela, e variam de acordo com o algoritmo escolhido como solução\cite{Kale:2010}. 
No exemplo do Page Rank, o número de páginas é enome, e elas são recolhidas de diversos servidores da Google; é uma questão fundamental escolher algoritmo paralelo com o melhor desempenho dentre as soluções possíveis.

\textit{Devido ao grande número de algoritmos de ordenação paralela e variadas arquiteturas paralelas, estudos experimentais assumem uma importância crescente na avaliação e seleçãoo de algoritmos apropriados para multiprocessadores.}

%3. porque map reduce

O MapReduce\cite{Dean:2008}  é um modelo de programação paralela criado pela Google para processamento e geração de grandes volumes de dados em \textit{clusters}. Esse modelo propõe simplificar a computação paralela e ser de fácil uso, abstraindo conceitos complexos da paralelização - como tolerância a falhas, distribuição de dados e balanço de carga - e utilizando duas funções principais: map e reduce. A complexidade do algoritmo paralelo não é vista pelo desenvolvedor, que pode se ocupar em desenvolver a solução proposta. O Hadoop \cite{Hadoop:2010} é uma das implementações do MapReduce, um \textit{framework} open source que provê o gerenciamento de computação distribuída. Foi desenvolvido por Doug Cutting em 2005, e é amplamente apoiada e utilizado pela Yahoo!.




\section{Objetivos}

comparar duas ou mais implementações de algoritmos paralelos de ordenação
pode citar o trabalho da Paula e dizer que é uma continuação dele

\begin{itemize}
\item Estudar a programação paralela aplicada à algoritmos de ordenação
\item Implementar uma solução no modelo MapReduce, com o software Hadoop
\item Comparar e analisar o desempenho das soluções com relação à quantidade de dados a serem ordenados, variabilidade dos dados de entrada e número máquinas utilizadas.
\end{itemize}

\section{Resultados esperados}

Com a realização do trabalho, busca-se ampliar e consolidar conhecimentos adquiridos na área de computação paralela, assim como
a capacidade de análise e desenvolvimento de algoritmos paralelos no modelo MapReduce. 


Ao final do trabalho, espera-se obter a implementação de algoritmos de ordenação paralela em ambiente Hadoop e uma análise comparativa de desempenho entre tais algoritmos.

\section{Metodologia}

O início do projeto será destinado ao estudo mais detalhado da computação paralela, em especial os algoritmos de ordenação paralela, dos fatores que influenciam o desempenho de tais algoritmos, o modelo MapReduce e a plataforma Hadoop. O passo seguinte é conhecer detalhadamente o algoritmo paralelo a ser implementado e definir as estratégias para sua implentação ambiente Hadoop. 
O algoritmo implementado deve ser cuidadosamente avaliado para verificar um funcionamento adequado com diferentes entradas e número de máquinas. 
Em seguida, serão realizados experimentos para testes de desempenho dos algoritmos com relação à quantidade de máquinas, quantidade de dados e conjunto de dados.  Os resultados obtidos serão analisados e permitirão comparar a performance dos algoritmos em cada situação. 


\section{Infraestrutura necessária}

A infra estrutura necessária ao desenvolvimento do projeto será fornecida pelo Laboratório de Redes e Sistemas (LABORES) do Departamento de Computação (DECOM). Esse laboratório possui um cluster formado por cinco máquinas Dell Optiplex 380, que serão utilizadas na realização dos testes dos algoritmos. Os algoritmos serão desenvolvidos em linguagem Java, de acordo com o modelo MapReduce, no ambiente Hadoop. 

O cluster a ser utilizado apresenta as seguintes características:
\begin{itemize}
\item 5 nodos
\item Processador Intel Core 2 Duo de 3.0 GHz
\item Disco rígido SATA de 500 GB 7200 RPM
\item Memória RAM de 4 GB
\item Placa de rede Gigabit Ethernet
\item Sistema operacional Linux Ubunbu 10.04 32 bits (kernel 2.6.\textbf{XX})
\item Sun Java JDK 1.6.0 19.0-b09 
\item Hadoop 0.20.2
\end{itemize}

\section{Cronograma de trabalho}

Na Tabela \ref{tab:cronograma} está descrito o cronograma esperado para o desenvolvimento do projeto. Cada atividade foi descrita para se adequar da melhor maneira ao tempo disponível do projeto, mas é possível que o cronograma seja refinado posteriormente, com a inclusão de novas atividades ou redistribuição das tarefas existentes. 

\begin{table}[h]

\renewcommand{\arraystretch}{1}
\setlength\tabcolsep{3pt}
\begin{center}
\begin{tabular}{| c | c | c | c | c | c | c | c | c | c | c |}
\hline

Atividade &Fev &Mar &Abr &Mai &Jun &Jul &Ago &Set &Out &Nov \\ \hline \hline
\ref{c1}   &$\bullet$ &$\bullet$ & & & & & & & & \\ \hline
\ref{c2}   & &$\bullet$ &$\bullet$ & & & & & & & \\ \hline
\ref{c3}   & & &$\bullet$ & & & & & & & \\ \hline
\ref{c4}   & & &$\bullet$ &$\bullet$ & & & & & & \\ \hline
\ref{c5}   & & & &$\bullet$ &$\bullet$ & & & & & \\ \hline
\ref{c6}   & & & & & &$\bullet$ &$\bullet$ & & & \\ \hline
\ref{c7}   & & & & & & & &$\bullet$ & & \\ \hline
\ref{c8}   & & & & & & & & &$\bullet$ & \\ \hline
\ref{c9}   & & & & & & & & & &$\bullet$ \\ \hline
\hline
\end{tabular}
\end{center}
\caption{Cronograma proposto para o projeto}
\label{tab:cronograma}
\end{table}

\begin{enumerate}
 \item \label{c1} Pesquisa bibliográfica sobre o tema do projeto e escrita da proposta
 \item \label{c2} Estudo mais detalhado dos algoritmos de ordenação paralela,  modelo MapReduce e Hadoop.
 \item \label{c3} Configuração do ambiente Hadoop no laboratório.
 \item \label{c4} Implementação e testes.
 \item \label{c5} Escrita, revisão e entrega do relatório. 
 \item \label{c6} Implementação e testes.
 \item \label{c7} Análise comparativa entre os resultados.
 \item \label{c8} Escrita e revisão do projeto final.
 \item \label{c9} Entrega e apresentação.
 \end{enumerate}
 
%
%Citações: \\
%\cite{Kale:2010} \\
%\cite{Manferdelli:2008} \\
%\cite{Dean:2008} \\
%\cite{Asanovic:2009}

% ********** REFERÊNCIAS **********
%\bibliographystyle{abnt-alf}	 % Existem ainda: abbrv, acm, alpha, amsalpha, amsplain
\bibliographystyle{abnt-alf}
\bibliography{bibliografia} % o nome do arquivo .bib com as referências

\end{document}
